\chapter{基于局部信息语义结构的三相记忆设计：从局部拓扑表征到全局语义绑定}
\label{chap:heterogeneous-three-phase-memory}

\begin{chapterlead}
前一章把情景记忆 $E$ 从三相闭环中单独取出，建立了它的数学表示 $\mathcal{E}(t)=(\mu_E(t),\mathcal{G}_E(t))$、
液态动力学与对应的工程实现，使“经历如何成为历史、历史如何成为结构”有了严格的落点。但回顾整个三相记忆
框架 $h$--$E$--$\Theta$ 到目前为止的推进，可以发现一个始终隐含却未曾明说的假设：进入记忆系统的信息，
无论来源是什么，都可以被同一种表示方式——token、embedding 或图节点——统一处理。本章要指出，这一假设
并不成立，而且恰恰是可持续运行的 Agent 记忆系统最容易出错的地方。

问题的根源在于：进入记忆系统的信息本身具有不同的信息几何。视觉场、触觉场、声音频谱、身体姿态与连续运动
轨迹，本质上具有高度局部性、连续性与拓扑结构；而人物关系、目标、概念、因果、社会关系与任务状态，则更多
表现为跨空间、跨时间的分布式关系结构。两类信息并不适合使用同一种表示方式，也不适合以同一种动力学写入
同一种存储。因此，一个完整的记忆架构必须同时回答两个问题：信息以什么结构存在？这些结构当前处于什么记忆
相态？前者是表示几何问题，后者是动力学相态问题，二者正交而不可互相替代。

本章的推进分为四个层次。第一个层次建立信息结构的分类学：定义局部拓扑结构 $L$ 与全局分布式关系结构 $D$，
论证二者不能统一压缩为单一数据结构，提出全章的核心原则——表示异构，身份统一；并给出跨结构共用的最小
语义协议 What/Where/When，它把异构表示锚定到同一个现实事件上。第二个层次建立双向通道：局部信息经事件化
$L\rightarrow D$ 进入高级认知，全局语义经投影 $D\rightarrow L$ 反向约束局部网络，二者合成
感知、理解、目标化感知的完整闭环。第三个层次把 $L/D$ 分解注入三相记忆：$h$、$E$、$\Theta$ 中的每一相
都分解为局部成分、分布式成分与二者绑定成分，三相记忆由此成为一个三行（相态）乘三列（结构）的矩阵式演化
系统；与之配套，本章给出认知绑定总线、认知事件与工作记忆异构黑板三项基础设施，并论证工作记忆保存的应当
是引用与激活摘要而非全量数据。第四个层次处理工程运行机制：局部高速闭环与全局低频闭环异步多速率运行，
事件检测器把连续世界切成离散认知事件流，绑定器负责“何时绑定、何时分离”，SPC、TCE、CCM 与 Boundary
组成认知控制面，Offloading 把部分长期绑定外化；在此基础上统一各相的写入机制与学习速度，并以模式分离与
模式补全保证绑定质量，最终把全部结构设计收束进统一的认知状态空间与动力学方程。

由此，三相记忆不再只是一个记忆分层模型，而成为一个明确承认表示异构、以绑定维持身份统一的持续学习动力
系统。这也为下一章重新打开学习问题准备了结构前提：当一次真实经历到来，系统需要改变的不只是某一处参数，
而是分布在异构结构之上的多个候选对象——究竟应当改变什么，正是学习调度要回答的问题。
\end{chapterlead}

\section{问题的提出：记忆的信息几何}

传统 Agent 记忆系统通常把记忆划分为短期记忆、长期记忆、向量数据库、情景记忆、语义记忆等不同存储模块。
这种设计在工程上具有直观性，却容易忽略一个更基础的问题：

\begin{statementbox}
进入记忆系统的信息，本身具有不同的信息几何。
\end{statementbox}

视觉场、触觉场、声音频谱、身体姿态以及连续运动轨迹，本质上具有高度局部性、连续性和拓扑结构；而人物关系、
目标、概念、因果、社会关系和任务状态，则更多表现为跨空间、跨时间的分布式关系结构。两类信息并不适合使用
同一种表示方式。

因此，一个可持续运行的 Agent 记忆系统不应试图将所有信息统一压缩为 token、embedding 或图节点，而应承认：

\begin{statementbox}
智能内部至少同时存在两种异构的信息结构。
\end{statementbox}

即局部拓扑结构 $L$ 与全局分布式关系结构 $D$。其中 $L$ 主要承载连续、高频、邻域相关的信息，而 $D$ 主要
承载抽象、关系化、跨时间和跨空间的信息。

三相记忆
\[
h(t)\leftrightarrow E\leftrightarrow\Theta
\]
则并不是与 $L$、$D$ 并列的第三种分类。它描述的是这些信息在智能系统中的动力学相态：$h(t)$ 是当前工作
认知面，$E$ 是情景经历轨迹，$\Theta$ 是跨生命周期的长期生成结构。

因此，一个更加完整的 Agent 记忆架构应当同时解决两个问题：

\begin{statementbox}
\centering
\adjustbox{max width=.96\linewidth}{\(\displaystyle
\text{信息以什么结构存在？}
\qquad
\text{这些结构当前处于什么记忆相态？}
\)}
\end{statementbox}

由此形成本章提出的工程模型：

\begin{statementbox}
\centering
\adjustbox{max width=.96\linewidth}{\(\displaystyle
\text{局部—分布式异构信息结构}
\times
\text{三相记忆动力学}
\)}
\end{statementbox}

前一因子由本章展开，后一因子由本书此前各章建立，二者是正交的两个维度。

\section{两类基本信息结构}

\subsection{局部拓扑信息结构}

定义局部结构 $L$ 为满足较强邻域保持性质的信息表示。若输入空间中的两个状态 $x_i,x_j$ 在物理、空间或连续
信号意义上接近：
\[
d(x_i,x_j)\rightarrow0,
\]
则其内部表示通常也保持较强邻近关系：
\[
d(z_i,z_j)\rightarrow0.
\]

这一类型典型包括：视觉空间、身体表面、声音频谱、触觉阵列、空间地图、局部运动轨迹以及连续控制状态。其
主要特征是：

\begin{statementbox}
\centering
\adjustbox{max width=.96\linewidth}{\(\displaystyle
\text{连续}+\text{高频}+\text{高带宽}+\text{局部相关}
\)}
\end{statementbox}

例如机器人视觉系统每一时刻可能处理 $10^5\sim10^7$ 级别的原始感知变量。这些变量之间存在大量局部空间
相关性，但其中绝大部分不需要直接进入高级认知。

因此，$L$ 不能被直接等价为“语义”。它首先是现实世界在 Agent 内部形成的局部结构场。

\subsection{全局分布式关系结构}

定义 $D$ 为以关系而不是物理邻域组织的信息结构。例如：
\[
\text{Alice}\xrightarrow{\mathrm{owns}}\text{Cup},
\qquad
\text{Cup}\xrightarrow{\mathrm{used\ for}}\text{Drink},
\qquad
\text{Alice}\xrightarrow{\mathrm{requested}}\text{Give(Cup)}.
\]

这些信息之间是否相关，并不由物理距离决定，而由关系函数
\[
R(x_i,x_j)
\]
决定。因此，$L$ 更接近 distance-dominated representation，而 $D$ 更接近 relation-dominated
representation。

典型的 $D$ 包括：概念、对象身份、人物关系、社会关系、任务图、因果结构、故事线、计划、自我以及长期世界
模型。其主要特征是：

\begin{statementbox}
\centering
\adjustbox{max width=.96\linewidth}{\(\displaystyle
\text{稀疏}+\text{关系化}+\text{跨模态}+\text{上下文相关}
\)}
\end{statementbox}

\section{表示异构，身份统一}

一个重要工程原则是：

\begin{statementbox}
\centering
\adjustbox{max width=.96\linewidth}{\(\displaystyle
L\neq D
\)}
\end{statementbox}

不能简单地使用一种统一的数据结构表示全部认知。如果将完整视觉场全部转换为语义节点，则 $L\rightarrow D$
会导致巨大的信息膨胀；相反，如果试图使用局部连续网络表达复杂社会关系，又会导致结构表达能力不足。因此
正确目标不是把二者压缩进第三种结构：
\[
L+D\rightarrow N,
\]
而是保持二者的异构表示，并建立它们之间的绑定：
\[
\boxed{
L
\overset{\mathrm{Binding}}{\rightleftarrows}
D
}
\]

也就是说：

\begin{statementbox}
局部和全局不需要具有同一种表示；真正需要统一的是它们对同一个现实对象、事件和状态的指向。
\end{statementbox}

这形成本章最核心的工程原则：

\begin{statementbox}
\centering
\adjustbox{max width=.96\linewidth}{\(\displaystyle
\text{Representation Heterogeneity}+\text{Identity Unity}
\)}
\end{statementbox}

即：表示异构，身份统一。

\section{What/Where/When：跨结构的共同语义协议}

如果 $L$ 和 $D$ 使用不同的信息几何，那么二者之间必须存在一个最低公共协议。本章认为
What/Where/When 正好构成这个协议。任何进入认知界面的现实事件都可以首先抽象为：
\[
e_i=(Q_i,X_i,\tau_i),
\]
其中 $Q=\mathrm{What}$ 表示当前认知对象或状态是什么；$X=\mathrm{Where}$ 表示它处于哪个有效状态
空间；$\tau=\mathrm{When}$ 表示其时间、顺序、阶段和变化位置。

需要特别强调：Where 不仅指三维空间位置。它也可以表示任务空间位置、社会关系位置、逻辑状态位置以及抽象
状态空间位置。同样，When 也不仅是时间戳，而包括事件顺序、行为阶段、因果时序以及生命周期位置。

因此，What/Where/When 实际上构成的是一种认知事件坐标（Cognitive Event Coordinate）：无论一个事件
的底层表示属于 $L$ 还是 $D$，它都必须能在这个坐标系中被定位，跨结构的绑定才有公共的参照系。

\section{双向通道：事件化与投影}

\subsection{局部信息的事件化}

局部信息进入高级认知之前，需要经过一个关键过程——事件化（Eventization）：
\[
\text{连续场}\rightarrow\text{对象}\rightarrow\text{事件}\rightarrow\text{关系}.
\]

可以写成：
\[
L_t
\xrightarrow{\mathcal O}
O_t
\xrightarrow{\mathcal E}
\mathrm{Event}_t
\xrightarrow{\mathcal R}
D_t
\]
其中 $\mathcal O$ 负责对象化，$\mathcal E$ 负责事件化，$\mathcal R$ 负责关系化。

例如原始视觉输入 $\mathrm{Image}_t$ 首先形成对象 $\mathrm{Cup}_{17}$，然后形成事件
$\mathrm{Cup}_{17}\ \mathrm{moved}$，最终形成关系 $\mathrm{Alice\ moved\ Cup}_{17}$。于是高维局部
状态被压缩成有限认知事件。因此，$L\rightarrow D$ 的本质不是“编码格式转换”，而是抽象（Abstraction）。

\subsection{分布式结构的反向投影约束}

如果认知流程只有 $L\rightarrow D$，那么系统仍然只是感知系统。真正的智能必须同时具有反向通道
$D\rightarrow L$。例如高级语义系统形成目标：
\[
\mathrm{Goal}=\text{Find the red cup},
\]
这个目标并不应该直接控制视觉网络的每个内部状态。它更合理的作用是产生注意掩码（AttentionMask）、
特征增益（FeatureGain）、兴趣区域（RegionOfInterest）或者检索查询（SearchQuery）。

因此：
\[
D\xrightarrow{\mathcal P}L,
\qquad
\mathcal P(D)=\{\mathrm{gain},\mathrm{mask},\mathrm{query},\mathrm{constraint},\mathrm{target}\},
\]
其中 $\mathcal P$ 是投影算子。也就是说：

\begin{statementbox}
全局语义通过改变局部网络的边界条件，而不是通过逐状态遥控，实现对局部网络的控制。
\end{statementbox}

因此真正完整的闭环是：
\[
\boxed{
L\rightarrow D\rightarrow L
}
\]
对应：
\[
\boxed{
\text{感知}\rightarrow\text{理解}\rightarrow\text{目标化感知}
}
\]
进一步则是 $D\rightarrow L_{\mathrm{motor}}\rightarrow\mathrm{Action}$，形成
“理解 $\rightarrow$ 局部控制 $\rightarrow$ 行动”。

\section{三相记忆的异构表示}

在此前的三相记忆模型中，$h(t)$ 表示当前工作认知面，$E$ 表示情景经历轨迹，$\Theta$ 表示跨生命周期的
长期生成结构。引入 $L/D$ 以后，三相中的每一相都不能继续被看成单一结构。

首先，工作记忆应当分解为：
\[
\boxed{
h=h_L\oplus h_D\oplus h_{LD}
}
\]
其中 $h_L$ 表示当前被激活的局部拓扑状态，$h_D$ 表示当前被激活的分布式关系状态，而 $h_{LD}$ 表示当前
正在维持的局部—语义绑定关系。仅仅写成 $h=h_L\oplus h_D$ 仍然不完整，因为真正重要的是
$L\leftrightarrow D$ 如何发生绑定。

同样，情景记忆不能仅仅表示为 $E_L+E_D$。因为一次经历最重要的信息之一恰恰是：哪个局部状态，在什么上下文
中，被理解成什么意义？因此：
\[
\boxed{
E=E_L\oplus E_D\oplus E_{LD}
}
\]
其中 $E_L$ 保存局部轨迹，例如空间轨迹、身体运动、视觉变化、声音序列；$E_D$ 保存对象身份、事件意义、
人物关系、目标状态和因果解释；而 $E_{LD}$ 保存局部结构与语义结构在具体经历中的绑定历史。

例如，局部身体轨迹
\[
L:\ \text{hand raised + lateral movement}
\]
被绑定为
\[
D:\ \text{Alice waved goodbye}.
\]
这条对应关系 $L\leftrightarrow D$ 本身就是情景记忆的重要组成部分。

长期学习以后，$E_L\rightarrow\Theta_L$ 形成稳定局部规律，$E_D\rightarrow\Theta_D$ 形成稳定语义和关系
结构，同时 $E_{LD}\rightarrow\Theta_{LD}$ 形成两种信息空间之间的长期映射规律。因此：
\[
\boxed{
\Theta=\Theta_L\oplus\Theta_D\oplus\Theta_{LD}
}
\]
其中 $\Theta_L$ 负责局部感知规律、运动规律、身体模型、空间模式以及局部技能；$\Theta_D$ 负责概念、
世界模型、因果结构、社会关系、自我结构以及长期目标体系；而 $\Theta_{LD}$ 则负责 grounding、
abstraction、projection，也就是回答“什么样的局部结构意味着什么”，以及“什么样的全局目标应当转换为
什么样的局部约束”。

三相记忆因此获得新的完整形式：
\[
\boxed{
h=(h_L,h_D,h_{LD}),\qquad
E=(E_L,E_D,E_{LD}),\qquad
\Theta=(\Theta_L,\Theta_D,\Theta_{LD})
}
\]
这意味着三相记忆不再只是 $h\rightarrow E\rightarrow\Theta$ 的线性沉降，而是一个矩阵式演化：
\[
\begin{bmatrix}h_L & h_D & h_{LD}\end{bmatrix}
\ \longrightarrow\
\begin{bmatrix}E_L & E_D & E_{LD}\end{bmatrix}
\ \longrightarrow\
\begin{bmatrix}\Theta_L & \Theta_D & \Theta_{LD}\end{bmatrix}
\]
其中三列分别表示局部结构、分布结构、跨结构绑定；三行分别表示当前状态、经历轨迹、稳定生成规律。

\section{绑定基础设施：认知绑定总线与认知事件}

工程上最重要的模块，因此不应只是普通消息总线，而应定义为认知绑定总线（Cognitive Binding Bus）。
它的职责不是传输所有数据，而是维护：
\[
\boxed{
\text{现实对象}\leftrightarrow\text{局部表示}\leftrightarrow\text{语义身份}\leftrightarrow\text{事件轨迹}
}
\]

最基本的实体结构可以定义为：

\begin{verbatim}
CognitiveEntity {
    entity_id

    local_ref
    semantic_ref

    current_event_id

    confidence
    salience

    self_relevance
    task_relevance
}
\end{verbatim}

其中 \verb|local_ref| 指向局部网络中的当前表示，例如视觉对象、声音对象、空间对象、身体节点；
\verb|semantic_ref| 则指向概念节点、人物节点、知识图节点、任务节点。因此 $\mathrm{entity\_id}$ 成为
$L$ 和 $D$ 之间的共同身份锚点。

事件层进一步使用认知事件（Cognitive Event）：

\begin{verbatim}
CognitiveEvent {
    event_id

    what
    where
    when

    participants[]

    local_refs[]
    semantic_refs[]

    context_id

    confidence
    novelty
    salience
    self_relevance
    task_relevance
}
\end{verbatim}

它本身不保存所有感知细节，而是作为跨空间索引（Cross-space Index）指向不同表示空间。因此认知事件
本质上是 What/Where/When 与绑定的合成，这就是局部信息进入三相记忆系统的最基本工程单元。

\section{工作记忆的黑板结构：引用而非复制}

因此 $h(t)$ 最适合实现为一个异构认知黑板（Heterogeneous Cognitive Blackboard），而不是统一 token
buffer。可以划分为：

\begin{verbatim}
h(t)

├── Local State Space
│   ├── vision
│   ├── audio
│   ├── spatial
│   ├── body
│   └── motor
│
├── Binding Space
│   ├── entities
│   ├── events
│   ├── What
│   ├── Where
│   └── When
│
└── Distributed State Space
    ├── semantic graph
    ├── goals
    ├── task graph
    ├── episodic context
    ├── causal structure
    ├── social relations
    └── self-related states
\end{verbatim}

这里 Local State Space 不需要把全部状态复制到 Blackboard，它只维护引用与激活摘要（reference +
active summary）；Distributed State Space 也不是完整长期知识库，它只维护当前真正参与认知的语义子图。
因此 $h(t)$ 仍然保持有限。

如果视觉特征、空间地图和完整语义图全部复制进 $h(t)$，工作记忆复杂度将迅速爆炸。因此 $h(t)$ 真正需要
保存的是激活引用（Active References），而不是全量数据（Full Data）。例如 $h_L$ 可以只保存
\verb|object_17|、位置与视觉特征引用，而不是保存完整图像；同样，$h_D$ 只保存当前激活语义节点，而不是
完整知识图。所以：

\begin{statementbox}
\centering
\adjustbox{max width=.96\linewidth}{\(\displaystyle
h(t)=\text{有限激活索引面}
\)}
\end{statementbox}

比“工作记忆存储区”更加准确地描述了 $h(t)$ 的工程实质。

\section{多速率异步双闭环}

局部网络通常具有更高更新频率，例如 $\omega_L\gg\omega_D$。局部系统可能持续进行视觉跟踪、身体控制、
触觉反馈、局部运动预测，而全局网络并不需要参与每一次局部状态变化。

因此 AgentOS 不应该采用“Perception $\rightarrow$ LLM $\rightarrow$ Reasoning $\rightarrow$ Action
$\rightarrow$ repeat”这样的统一串行主循环，更合理的是异步多速率架构（Asynchronous Multi-rate
Architecture）。局部闭环
\[
L_t\rightarrow\mathrm{Action}\rightarrow\mathrm{World}\rightarrow L_{t+\delta}
\]
高速持续运行；高级认知闭环
\[
D_t\rightarrow\mathrm{Goal}\rightarrow\mathrm{Event}\rightarrow D_{t+\Delta}
\]
以较低频率或事件驱动方式运行，满足 $\delta\ll\Delta$。

因此整个系统天然形成两个基本闭环。局部高速闭环
\[
\boxed{
L\rightarrow\mathrm{Action}\rightarrow\mathrm{World}\rightarrow L
}
\]
主要解决稳定控制、连续感知、实时反应；全局低频闭环
\[
\boxed{
D\rightarrow\mathrm{Goal}\rightarrow\mathrm{Interpretation}\rightarrow D
}
\]
主要解决理解、规划、语义推理、目标调整。两者之间通过抽象通路
$L\xrightarrow{\mathrm{Abstraction}}D$ 和投影通路 $D\xrightarrow{\mathrm{Projection}}L$ 实现耦合。
因此 AgentOS 不是单循环系统，而是：

\begin{statementbox}
\centering
\adjustbox{max width=.96\linewidth}{\(\displaystyle
\text{Fast Local Loop}\rightleftarrows\text{Slow Distributed Loop}
\)}
\end{statementbox}

\section{事件检测与绑定计算}

\subsection{事件检测器：两层之间的频率桥}

如果所有局部变化都进入全局语义系统，那么 $D$ 会被大量无意义变化淹没。因此必须引入事件检测器
（Event Detector），其任务是判断 $\Delta L$ 是否值得进入高级认知。可以定义事件显著度：
\[
\boxed{
S_{\mathrm{event}}
=
w_1N+w_2G+w_3R_\Psi+w_4R_{\mathrm{task}}+w_5U
}
\]
其中 $N$ 表示新颖度，$G$ 表示局部变化幅度，$R_\Psi$ 表示自我相关度，$R_{\mathrm{task}}$ 表示任务
相关度，$U$ 表示不确定性。当
\[
S_{\mathrm{event}}>\theta_e
\]
时才产生认知事件并进入 $h_D$。这使局部高速连续世界被转换成离散的认知事件流。

\subsection{绑定器：异构认知真正的核心}

事件化之后仍然存在另一个问题：一个局部状态究竟应该与哪个语义节点绑定？因此需要独立的绑定器
（Binder）。其输入可以包括 What、Where、When、Context、Goal 与 $\Psi$。Binder 计算绑定概率：
\[
P_{\mathrm{bind}}(i,j)
=
\sigma\!\left(
\alpha S_{\mathrm{semantic}}
+\beta S_{\mathrm{spatial}}
+\gamma S_{\mathrm{temporal}}
+\eta S_{\mathrm{context}}
+\mu S_{\mathrm{goal}}
+\nu S_\Psi
\right),
\]
然后决定 $l_i\leftrightarrow d_j$ 是否成为稳定当前绑定。

高质量记忆系统的重要能力不仅是建立联系，还包括分离（separation）。如果任何共现信息都被绑定，关系图
将发生不可控扩散。因此 $P_{\mathrm{bind}}$ 还需要包含冲突项：
\[
P_{\mathrm{bind}}(i,j)
=
\sigma\!\left(
S_{\mathrm{support}}-\lambda C_{\mathrm{conflict}}
\right),
\]
其中 $C_{\mathrm{conflict}}$ 表示此次绑定与已有身份、因果关系和上下文结构发生冲突的程度。因此
Binder 的真正任务是：

\begin{statementbox}
\centering
\adjustbox{max width=.96\linewidth}{\(\displaystyle
\text{Bind when needed, Separate when needed}
\)}
\end{statementbox}

\section{控制回路：SPC、TCE 与 CCM}

Binder 不能自行无限建立关系，因为 $h(t)$ 容量有限。设 $R_L$ 为局部计算资源，$R_D$ 为分布式语义
计算资源，$R_B$ 为绑定计算资源，则：
\[
\boxed{
R_L+R_D+R_B\leq C_h
}
\]
CCM 的任务就是求解资源分配 $(R_L^*,R_D^*,R_B^*)$ 使当前任务效用最大：
\[
\mathbf R^*
=
\arg\max_{\mathbf R}
U(\mathrm{Task},\mathrm{World},h,E,\Theta,\Psi),
\qquad
\sum_iR_i\leq C_h .
\]

不同任务对应不同资源分配。例如机器人正在高速避障时 $R_L\gg R_D$，局部空间和身体控制优先；停下来制定
维修计划时 $R_D\gg R_L$，语义、因果和规划占据认知资源；如果机器人看到一个异常零件，则 $R_B\uparrow$，
因为系统需要重新识别“局部结构 $\leftrightarrow$ 已有语义”之间的关系。因此智能状态的变化可以表现为
\[
\boxed{
L\rightarrow B\rightarrow D\rightarrow B\rightarrow L
}
\]
的资源迁移。

TCE 不负责具体分配，而负责全局动力学。例如面对高度不确定环境时 $T\uparrow$，意味着更宽的搜索范围、
更频繁的假设切换、更高的新信息接纳能力；进入稳定执行阶段以后 $T\downarrow$，意味着更集中、更确定、
更稳定的局部控制。因此：
\[
\boxed{
\mathrm{TCE}\rightarrow\mathrm{Global\ Regime},
\qquad
\mathrm{CCM}\rightarrow\mathrm{Resource\ Allocation}
}
\]

CCM 和 TCE 如果没有对当前认知状态的理解，就无法正确控制。因此 SPC 沿自我 $\Psi$ 方向对当前认知状态
进行压缩：
\[
h(t)\xrightarrow{\mathrm{SPC}}\hat h_\Psi(t),
\]
得到“当前‘我’正在感知、理解和执行什么”的状态估计。这使整个控制闭环形成：
\[
h\rightarrow\mathrm{SPC}\rightarrow\mathrm{TCE}\rightarrow\mathrm{CCM}\rightarrow\mathrm{Binder}\rightarrow h'.
\]

\section{跨结构长期映射与具身语义}

一次局部信息进入 Agent，可以经历如下三相流转：
\[
L_{\mathrm{raw}}\ \longrightarrow\ h_L.
\]
如果被 Event Detector 判断具有认知意义，则
\[
h_L\rightarrow h_{LD}\rightarrow h_D.
\]
之后形成情景：
\[
(h_L,h_{LD},h_D)\rightarrow(E_L,E_{LD},E_D).
\]
经过多次类似经历：
\[
E_L\rightarrow\Theta_L,
\qquad
E_D\rightarrow\Theta_D,
\qquad
E_{LD}\rightarrow\Theta_{LD}.
\]
于是局部经验逐渐转变成稳定能力。

$\Theta_{LD}$ 是具身语义形成的核心。如果一个 Agent 只有 $\Theta_D$，它可以拥有大量语言知识，却可能
无法将这些知识可靠地绑定到真实世界；如果只有 $\Theta_L$，它可以拥有非常强的视觉和运动技能，却未必
真正拥有抽象理解。真正的具身智能需要 $\Theta_{LD}$ 建立 perception $\leftrightarrow$ meaning 与
meaning $\leftrightarrow$ action 之间长期稳定的映射。从这个角度看，$\Theta_{LD}$ 实际上就是扎根智能
（Grounded Intelligence）长期形成的结构基础。

工程上，长期结构层可以直接分成三类模型：

\begin{verbatim}
Theta

├── Local Models
│   ├── perception model
│   ├── spatial model
│   ├── body model
│   └── skill model
│
├── Distributed Models
│   ├── semantic graph
│   ├── causal model
│   ├── world model
│   ├── social model
│   └── self model
│
└── Cross-space Models
    ├── grounding model
    ├── object identity model
    ├── event abstraction model
    ├── goal-to-action model
    └── semantic-to-attention model
\end{verbatim}

其中 Cross-space Models 即 $\Theta_{LD}$。

\section{生物启发：海马式绑定与前额叶式控制}

人脑中的海马系统可以给这一架构提供重要的功能启发，但工程设计不需要复制其生理结构。可以抽象出四个能力：
第一是索引（Indexing），保存事件中分散表示的访问索引；第二是绑定（Binding），把 What、Where、
When、对象身份和上下文绑定起来；第三是模式分离（Pattern Separation），避免相似事件被错误融合；
第四是模式补全（Pattern Completion），在获得部分线索时重建整个事件结构。因此一个工程 Binder 可以
设计为：

\begin{statementbox}
\centering
\adjustbox{max width=.96\linewidth}{\(\displaystyle
\text{Index}+\text{Bind}+\text{Separate}+\text{Complete}
\)}
\end{statementbox}

而不必复制真实海马神经回路。

同样，前额叶提供的不是一个“中央数据库”模板，真正有价值的是控制面（Control Plane）概念：系统需要
一个控制平面决定当前目标是什么、哪些信息值得进入 $h(t)$、哪些 Binder 应被激活、哪些已有记忆应该被
召回、哪些关联应该被抑制、当前应该继续搜索还是开始执行。因此 AgentOS 中 CCM+TCE+SPC 共同构成
认知控制面（Cognitive Control Plane），而 $L$+$D$+Binder 构成认知数据面（Cognitive Data Plane）。
这是一个非常重要的工程分层。

\section{控制面与数据面的分层}

因此整个 AgentOS 可以进一步分成两个大层。数据面
\[
\boxed{
L+D+\mathrm{Binding}+h/E/\Theta
}
\]
负责实际认知状态的产生、流转和记忆；控制面
\[
\boxed{
\mathrm{SPC}+\mathrm{TCE}+\mathrm{CCM}+\mathrm{Boundary}
}
\]
负责观测、调节、资源分配、输入输出边界管理。这类似计算系统中 Data Plane 与 Control Plane 的分离。

局部信息通常拥有巨大带宽，所以真正进入 $h_L$ 之前还必须经过边界（Boundary）过滤。可以定义：
\[
\mathrm{Input}_{\mathrm{accepted}}
=
B(\mathrm{Input}_{\mathrm{world}},\mathrm{Goal},\Psi,C_h).
\]
Boundary 根据当前任务、资源余量、自我相关性、不确定性以及环境风险，动态调整输入阈值。因此 Boundary
解决的是“多少世界应该进入 Agent”，CCM 解决“进入以后资源怎么分”，Binder 解决“不同结构之间怎么绑定”。

某些结构不需要持续保存在 $\Theta$ 内部。例如复杂地图、工具说明、大量历史文件、完整视觉记录，可以外化
到外部环境 $M_e$ 中。于是内部只需要保持“如何访问、何时访问、访问后如何重新绑定”，即 $\Theta_{LD}$
甚至可以保存 semantic node $\leftrightarrow$ external resource 之间的映射。因此 Offloading 不只是
减小存储量，它同样减少内部持续绑定复杂度。

\section{运行时架构与完整认知周期}

一个实际 AgentOS 可以采用如下运行时结构：

\begin{verbatim}
World / Body
    |
    v
Local Perception & Control Networks
    |
    |----------> Fast Local Action Loop
    |
    v
Event Detector
    |
    v
Object / Event Binder
    |
    v
Cognitive Binding Bus
    |
    v
Heterogeneous h(t)
 +---------+------------+---------+
 | Local   | Binding    | Global  |
 | State   | State      | State   |
 +---------+------------+---------+
    |
    |----> SPC
    |        |
    |        v
    |       TCE
    |        |
    |        v
    |       CCM
    |        |
    |        └────> allocation / suppression
    |
    v
Distributed Semantic / Goal / World Models
    |
    |----> semantic constraint
    |
    v
Projection Layer
    |
    v
Local Networks
\end{verbatim}

与此同时，$h\rightarrow E\rightarrow\Theta$ 持续运行在后台结构更新路径中。

一个完整认知周期可以描述为十三个步骤：第一步，局部网络持续运行 $L_t$；第二步，Boundary 过滤输入；
第三步，Event Detector 检测显著变化 $\Delta L_t$；第四步，对局部状态进行对象化和事件化
$L\rightarrow\mathrm{What/Where/When}$；第五步，Binder 将当前局部对象与已有语义节点绑定
$L\leftrightarrow D$；第六步，绑定后的状态进入 $h(t)$；第七步，SPC 获取当前认知状态摘要；第八步，
TCE 决定全局动力学模式；第九步，CCM 决定 $R_L,R_D,R_B$；第十步，全局语义系统产生新的目标或解释；
第十一步，通过 Projection $D\rightarrow L$ 改变局部注意、查询和控制边界；第十二步，具有长期价值的
状态进入 $E$；第十三步，重复稳定结构进一步形成 $\Theta$。这构成完整闭环。

\section{写入机制与多时间尺度学习}

记忆写入不应该是统一操作。不同三相、不同结构，其写入机制应完全不同。对于 $h_L$，更新应该是高频覆盖式；
对于 $h_D$，应该是事件驱动式；对于 $E$，应该以认知事件为写入基本单位；对于 $\Theta$，应该以稳定不变量
（stable invariant）为写入单位；而 $\Theta_{LD}$ 只有在 $L\leftrightarrow D$ 绑定关系反复稳定出现以后才
应该更新。这能够避免一次偶然经历直接污染长期结构。

可以定义三种学习速度：
\[
\boxed{
\eta_h\gg\eta_E\gg\eta_\Theta
}
\]
其中 $\eta_h$ 表示当前状态更新率，$\eta_E$ 表示情景结构更新率，$\eta_\Theta$ 表示长期结构更新率。跨结构
绑定 $\Theta_{LD}$ 通常还应采用更加保守的长期更新：
\[
\eta_{\Theta_{LD}}\leq\eta_\Theta,
\]
因为错误 grounding 往往比单一知识错误更加危险。

\section{模式分离、模式补全与记忆召回}

系统一个重要失败模式是错误绑定（False Binding）。例如两个长相相似的人被绑定到同一个身份，或者两个具有
相似局部模式的事件被错误理解成同一种语义。因此 Binder 必须同时维护相似度（Similarity）与差异度
（Difference），可以定义：
\[
\mathrm{Score}_{\mathrm{bind}}
=
\mathrm{Similarity}
-\lambda\,\mathrm{DifferenceCriticality}.
\]
如果 DifferenceCriticality 足够高，即使总体相似度很高也不应绑定。这使模式分离成为三相记忆工程中的
基础机制。

反过来，当当前局部输入 $L'$ 只包含过去事件的一部分时：
\[
L'\subset L_i,
\]
Binder 可以通过 $\Theta_{LD}$ 找到对应 $D_i$，随后召回 $E_i$，形成：
\[
L'\rightarrow D_i\rightarrow E_i.
\]
这就是模式补全（Pattern Completion）。因此记忆召回不是 $\mathrm{query}\rightarrow\mathrm{document}$，
而是：
\[
\boxed{
\text{partial state}\rightarrow\text{binding hypothesis}\rightarrow\text{episodic reconstruction}
}
\]

\section{统一认知状态空间与最终动力学}

局部—分布式结构与 What/Where/When 的关系，是两个不同维度：前者是信息结构几何，后者是事件功能坐标。
所以：
\[
\mathrm{What}=\mathrm{What}_L\oplus\mathrm{What}_D,
\qquad
\mathrm{Where}=\mathrm{Where}_L\oplus\mathrm{Where}_D,
\qquad
\mathrm{When}=\mathrm{When}_L\oplus\mathrm{When}_D.
\]
例如 $\mathrm{Where}_L$ 是物体物理位置，而 $\mathrm{Where}_D$ 可能表示其在任务流程中的位置。因此
$h(t)$ 可以进一步表示为：
\[
h(t)
=
\begin{bmatrix}
h_{\mathrm{What},L} & h_{\mathrm{What},D}\\
h_{\mathrm{Where},L} & h_{\mathrm{Where},D}\\
h_{\mathrm{When},L} & h_{\mathrm{When},D}
\end{bmatrix}
+h_{\mathrm{Binding}}.
\]
这是运行时认知状态更加完整的内部结构。

最终，局部信息语义结构设计可以被放入此前建立的认知状态空间：
\[
\boxed{
\mathcal M
=
\mathrm{Phase}\times\mathrm{SemanticCoordinate}\times\mathrm{RepresentationGeometry}
}
\]
其中 $\mathrm{Phase}=\{h,E,\Theta\}$，$\mathrm{SemanticCoordinate}=\{\mathrm{What},\mathrm{Where},
\mathrm{When}\}$，$\mathrm{RepresentationGeometry}=\{L,D,LD\}$。这里 $LD$ 不是第三种基础表示几何，
而是前两种结构之间的绑定关系。因此更严格地说：
\[
\boxed{
\mathcal M
=
\bigl(
\{h,E,\Theta\}
\times
\{\mathrm{What},\mathrm{Where},\mathrm{When}\}
\times
\{L,D\}
\bigr)
+\mathrm{Binding}
}
\]

完整系统的动力学模型可以写成。局部状态：
\[
\boxed{
h_L(t+\delta)
=
F_L(\mathrm{World},\mathrm{Body},h_L,\Theta_L,P_{D\rightarrow L})
}
\]
全局认知状态：
\[
\boxed{
h_D(t+\Delta)
=
F_D(h_D,E_D,\Theta_D,B_{L\rightarrow D},\mathrm{Goal},\Psi)
}
\]
绑定状态：
\[
\boxed{
h_{LD}
=
F_B(h_L,h_D,E_{LD},\Theta_{LD},\mathrm{Context})
}
\]
其中 $B_{L\rightarrow D}$ 表示抽象通路，而 $P_{D\rightarrow L}$ 表示投影通路。资源满足
$R_L+R_D+R_B\leq C_h$，并由 CCM 动态调节。

\section{本章小结：最终工程定义}

本章提出的“基于局部信息语义结构的三相记忆”可以最终定义为：

\begin{statementbox}
一种在承认局部拓扑信息与全局分布式语义信息具有异构表示几何的基础上，通过 What/Where/When 事件坐标
和 Cognitive Binding Bus 建立跨空间身份绑定，并将局部状态、语义状态及其绑定关系分别组织于工作态
$h(t)$、情景态 $E$ 与结构态 $\Theta$ 中，再由 SPC、TCE、CCM、Boundary 与 Offloading 对其进行观测、
调控、资源分配和外化管理的多频闭环记忆架构。
\end{statementbox}

其核心不在于建立更多记忆数据库，而在于解决四个基本问题：

\begin{statementbox}
什么信息应该进入认知？局部状态意味着什么？当前意义应该如何重新约束局部感知和行动？这种局部—语义绑定
何时应该成为长期结构？
\end{statementbox}

因此三相记忆的真正工程对象最终不再只是 Information，而是：
\[
\boxed{
\mathrm{State}+\mathrm{Relation}+\mathrm{Binding}+\mathrm{Dynamics}
}
\]
其中 $h$ 保存当前异构状态及其临时绑定，$E$ 保存这些绑定随经历形成的轨迹，$\Theta$ 保存局部规律、全局
规律以及跨表示空间的长期转换规律 $\Theta_{LD}$。

因此可以将整个架构最终浓缩为：
\[
\boxed{
L
\xrightarrow{\mathrm{Eventization}}
h_{LD}
\xrightarrow{\mathrm{Abstraction}}
D
\xrightarrow{\mathrm{Projection}}
L
}
\]
同时：
\[
\boxed{
h_{LD}\rightarrow E_{LD}\rightarrow\Theta_{LD}
}
\]
前一条构成实时认知闭环，后一条构成长期学习闭环，二者共同形成“感知—语义—行动”与“工作态—经历态—
结构态”的统一系统。

最终，Agent 并不是因为拥有一个巨大的统一模型而获得智能，而是在有限认知资源下，让高速局部网络、低频
全局语义网络以及两者之间的绑定机制，以不同频率、不同信息几何持续协同运行。这才是基于局部信息语义结构
的三相记忆系统真正的工程基础，也是下一章重新讨论学习调度时的结构前提。
